\chapter*{Summary of the investigation}
\addcontentsline{toc}{chapter}{Summary of the investigation}
\markboth{Summary of the investigation}{}

\section*{The concern}

The production shelf-life specialists reported internal test scores in the
range $\rsq \approx 0.92$--$0.97$ for the general-shelf-life task. A score
of that magnitude on a biological prediction problem is unusual enough to
warrant an explanation, and the question that was put to this project was
whether those scores reflected genuine predictive skill, classical
overfitting, or an evaluation procedure that was not sufficiently
independent of the training data.

That question is legitimate and this report treats it as such. The
investigation was not conducted to defend the original numbers. It was
conducted to determine what they measure.

\section*{What was done}

Six diagnostic experiments were executed against the production dataset,
covering model complexity reduction, a classification reformulation,
$K$-fold cross-validation at two values of $K$, category aggregation,
validation-set removal, and a model zoo spanning seventeen
configurations from a constant-output baseline to the production gradient
ensembles. The original production training and evaluation code was
reproduced from source rather than described from documentation, and the
structure of the synthetic corpus was audited directly.

\section*{What was found}

\begin{findingbox}[title=The central result]
Under the original evaluation protocol, \textbf{100\% of test rows in
every one of the six specialists originate from a synthetic generation
rule that is also present in that specialist's training set.} The
protocol correctly prevents a \emph{context} from appearing on both sides
of the split --- context overlap between train and test is exactly zero in
all six specialists --- but it places no constraint at all on the
generation rule. The held-out rows are therefore new samples of
relationships the model has already been trained on.
\end{findingbox}

This single structural fact explains the pattern of results better than
any model-side explanation, and the diagnostic experiments confirm it by
elimination:

\begin{itemize}[leftmargin=1.3em, itemsep=2pt]
\item Reducing model capacity does not reduce the score. A ridge
  regression --- a linear model with no interaction terms --- reaches
  $\rsq = 0.910$ on the soft general-shelf-life specialist where the
  production LightGBM reaches $0.971$. A problem that a linear model
  solves to within six points of a 2000-tree gradient ensemble is not a
  problem whose high scores are produced by model complexity.
\item Increasing $K$ in cross-validation does not reduce the score,
  because $K$-fold with the same grouping variable reshuffles the same
  already-seen rules.
\item Removing the validation set changes test $\rsq$ by a mean of
  $+0.0003$ across all model--specialist pairs: no material effect in
  either direction.
\item Aggregating all three cheese categories into a single training set
  does not reduce the score either.
\end{itemize}

\section*{The corrected evaluation}

The protocol was replaced with \emph{leave-one-generation-rule-out}
cross-validation: each fold removes every row belonging to one
\texttt{source\_rule\_id} from training and predicts that rule's rows
with a model that has never seen it. All seventeen configurations were
re-evaluated under this protocol across all six specialists.

\begin{findingbox}[title=Effect of the protocol change]
The three production tree ensembles retain a median fold $\rsq$ in the
region of $0.70$--$0.87$, substantially below their context-holdout
scores but clearly above a no-skill baseline. Several simpler models that
looked competitive under the original protocol collapse to strongly
negative $\rsq$ once the generation rule is withheld --- ridge regression
falls from a mean context-holdout $\rsq$ near $0.90$ to a strongly
negative leave-rule-out mean. The evaluation change, not any change to
the models, produces the entire difference.
\end{findingbox}

\section*{What this does and does not establish}

The revised numbers are an honest measure of one specific ability: the
ability to predict rows generated by a \emph{synthetic rule} the model has
not seen, using other synthetic rules as training data. They are not a
measure of real-world predictive accuracy on physical cheese, and this
report is explicit throughout about the difference. Chapter~\ref{ch:limits}
enumerates the dependencies that survive the corrected protocol, and
Chapter~\ref{ch:external} examines the external literature evaluation
that would be required to make a real-world claim.

\begin{limitbox}[title=Principal evidence limitation]
The synthetic data-generation source code is not present in this
repository; only its output is. Every statement in this report about
generation-rule behaviour is therefore inferred statistically from the
generated data, and is labelled as such. Separately, no dataset file in
this repository at any point in its recoverable history contains a
non-synthetic \texttt{data\_origin} tag, so the original experimental
observations that must have calibrated the generation rules cannot be
counted, characterised, or audited from here. Both limitations are
documented in full in \Cref{sec:limit-evidence}.
\end{limitbox}
